% Generated from locale/tr/content/first-order-logic/beyond/other-logics.tex; do not hand-edit.


\olfileid[tr]{fol}{byd}{oth}

\olsection{Diğer Mantıklar}

Şimdiye kadar anlamış olabileceğiniz gibi yeni bir mantık tasarlamak zor
değildir. Siz de kendi sözdiziminizi oluşturabilir, bir türetim sistemi
kurabilir ve ona uygun bir semantik geliştirebilirsiniz. !!{derivation}
sisteminin semantiğe göre tam olmasını istiyorsanız biraz ustalık göstermeniz,
mantığınızın gerçekten ilgi çekici olduğuna dünyayı ikna etmek için de emek
harcamanız gerekebilir. Buna karşılık mantığınızın matematiksel ve hesaplamalı
özelliklerini araştırarak saatlerce temiz ve keyifli vakit geçirebilirsiniz.

Son onyıllar biçimsel mantıklarda gerçek bir patlamaya tanık oldu. Bulanık
mantık, belirsiz özellikler hakkındaki akıl yürütmeyi modellemek üzere
tasarlanmıştır. Olasılıksal mantık, belirsizlik koşullarındaki akıl yürütmeyi
modellemeyi amaçlar. Varsayılan mantıklar ile monoton olmayan mantıklar,
çürütülebilir akıl yürütme biçimlerini, yani yeni bilgiler karşısında daha
sonra geri alınabilen ``makul'' çıkarımları modellemek üzere tasarlanmıştır.
Bilgi hakkındaki akıl yürütmeyi modelleyen epistemik mantıklar; nedensel
ilişkiler hakkındaki akıl yürütmeyi modelleyen nedensel mantıklar; hatta ahlaki
ve etik yükümlülükler hakkındaki akıl yürütmeyi modelleyen ``deontik''
mantıklar vardır. Bu sistemleri ortaya koymanın başlıca güdüsünün felsefi,
matematiksel ya da hesaplamalı olmasına bağlı olarak bu tür varlıkların
matematiksel mantık, felsefi mantık, yapay zekâ, bilişsel bilim ya da başka bir
başlık altında incelendiğini görebilirsiniz.

Liste uzayıp gider; olanaklar sonsuz görünür. Leibniz'ın bütün insan aklını
hesaplamaya indirgeme düşünün hiçbir zaman gerçekleştiremeyebiliriz---ama bu,
denememize engel olamaz.

